% sup2_07_cut: Supplement part 2: degrees of a graph after a line is cut (the degree facts used by the regrouping
% of Ba3 p. 432, which is proved in the main paper, Step 3(D_3)). Input from sup2_00_intro.tex's section.

\subsection{Degrees after a line is cut}\label{sup2:ssec:cut}

\cite[p.~432]{Ba3}: ``if we replace a line in an arbitrary graph $G$ by a pair of external fields, then the degree of the new graph is greater or equal to the degree of $G$, thus a convergent graph is transformed into a convergent one. From our analysis in previous chapters it follows that most of the divergent graphs are transformed into convergent ones, with the possible exception of graphs with one external vector field leg and vacuum graphs.'' The regrouping into classes which this sentence prepares is proved in the main paper; the three degree facts it uses are proved here. Throughout, a cut set $S$ is a set of lines of $G$, and $G|_S$ denotes $G$ with every line of $S$ replaced by its two legs, which become external fields (the kernel of the line, an external field by \cref{sup2:prop:S9,sup2:prop:S9vec}). Degrees are those of \cite[(2.1)--(2.3)]{Ba3}, $d=3$ where stated.

\begin{lemma}[degrees are intrinsic]\label{sup2:lem:cut1}
\textup{(a)} The connected subgraphs of $G|_S$ are exactly the connected subgraphs $H$ of $G$ with no line in $S$, and $D(H)$ is the same whether computed in $G|_S$ or in $G$. \textup{(b)} Hence, if every connected subgraph of $G$ has positive degree (the hypothesis of \cite[Prop.~2.1]{Ba3}, with or without the exception \cite[(2.4)]{Ba3}), so does every connected subgraph of $G|_S$; and the divergent connected subgraphs of $G|_S$ are exactly the divergent connected subgraphs of $G$ which avoid $S$.
\end{lemma}

\begin{proof}
A subgraph is given by its lines and their end vertices. \cite[(2.1)]{Ba3} for $D_H(v)$ counts only legs on lines of $H$, and a leg on a line of $S$ is on no line of $H$, both in $G$ and in $G|_S$.
\end{proof}

\begin{lemma}[the degree of the whole graph]\label{sup2:lem:cut2}
Apply \cite[(2.2)]{Ba3} verbatim to $G|_S$, with one $-d$ even if $G|_S$ is disconnected. Then $D(G|_S)-D(G)=\sum_{l\in S}\big[\sum_{\text{ends of }l}c_{\mathrm{end}}+\delta_l\big]$, with $c_{\mathrm{end}}=\frac{d-2}2$ for an ordinary leg, $c_{\mathrm{end}}=\frac{d-4}2$ for a vector leg of a vertex (1.14)/(1.15), and $\delta_l$ the number of differentiations acting on $l$. So the printed sentence ``the degree of the new graph is greater or equal to the degree of $G$'' holds whenever no line of $S$ ends at a vector leg of a (1.14)/(1.15) vertex, and at $d=3$ the degree drops by $\frac12$ per such end (two (1.14) vertices with $n=2$, $\phi'$ external, joined by their two vector lines have $D=5$; cutting one line gives $D=4$). What the regrouping uses is \cref{sup2:lem:cut1}, together with \cref{sup2:lem:S7cor}(a), which keeps every component containing a (1.13)--(1.15) vertex at $D\ge1$.
\end{lemma}

\begin{proof}
By \cref{sup2:lem:S7vert}(b), an ordinary leg leaving a line changes $D_G(v)$ by $+\frac{d-2}2$, a (1.14)/(1.15) vector leg by $-\frac{4-d}2$, and a no longer counted differentiation by $+1$. In the example, $D_G=3+2\cdot\frac12=4$ per vertex before and $3+\frac12$ after; with $-3$ once, $5$ and $4$.
\end{proof}

\begin{lemma}[cutting a divergent graph; $d=3$]\label{sup2:lem:cut3}
Let $\gamma$ be connected without (1.13)--(1.15) vertices, with $E(\gamma)\ge2$, and let $S'$ be a set of $t\ge1$ lines of $\gamma$ whose removal leaves $\gamma|_{S'}$ connected (for $t=1$: a line which is not a bridge; every line, if $\gamma$ is one-particle-irreducible) and such that $\gamma|_{S'}$ still has a line. Then
\[
D(\gamma|_{S'})=D(\gamma)+t+(\text{differentiations on }S')\ \ge\ \max\{0,D(\gamma)+t\},
\]
and $D(\gamma|_{S'})=0$ only if $t=1$ and $\gamma|_{S'}$ is the graph \cite[(2.4)]{Ba3}. If $E(\gamma)\ge3$, then $D(\gamma|_{S'})\ge1$.
\end{lemma}

\begin{proof}
\emph{Exact change.} By \cref{sup2:lem:S7id} for the connected $\gamma|_{S'}$: $\ell$ drops by $t$, $E$ rises by $2t$, $\Delta_{\mathrm{ext}}$ rises by the differentiations on $S'$, and $\kappa$ is unchanged. \emph{Lower bound.} Since the removal of $S'$ leaves a connected graph, $\gamma$ has a cycle, $\ell(\gamma)\ge1$ and $D(\gamma)\ge\ell+E-4\ge-1$; hence $D(\gamma|_{S'})\ge0$, with equality only if $t=1$, $D(\gamma)=-1$ and $S'$ carries no differentiation. \emph{Equality.} Then $\gamma|_{S'}$ has $\ell=0$, $E=4$, all $\kappa=0$ and $\Delta_{\mathrm{ext}}=0$, so it is \cite[(2.4)]{Ba3} by \cref{sup2:cor:S7class}, which requires a line; this is where the last hypothesis enters (without it, a self-loop cut of a one-vertex $\gamma$ leaves a bare vertex with four or two external legs and $D=0$: the \cite[(1.6)]{Ba3} tadpole and the \cite[(1.10)]{Ba3} tadpoles with $n+n'=2$). \emph{$E\ge3$.} $D(\gamma)\ge1+3-4=0$, so $D(\gamma|_{S'})\ge1$.
\end{proof}

The cut of \cite[p.~432]{Ba3} removes only lines joining distinct vertices whose localisations are at distance $\ge1$, never a self-loop, and a connected $\gamma|_{S'}$ with at least two vertices has a line; so the last hypothesis of \cref{sup2:lem:cut3} holds at every use. The outcome \cite[(2.4)]{Ba3} arises from the $e^2$ self-energy with its photon line cut and from the one-loop vacuum polarisation with one scalar line cut; it needs no partner, since the summation by parts of \cref{sup1:ssec:S4} treats it. The other two possible outcomes of cutting a divergent subgraph, a vacuum graph and a one-vector-leg graph, are Ba{\l}aban's exception list: vacuum classes always receive their counterterms \cite[p.~432]{Ba3}, and the one-vector-leg loop through a finite-counterterm vertex vanishes after the gauge transformation by \cref{sup2:lem:S7exc}(a) (its coefficient is an even polynomial in $q$). The regrouping itself, with these inputs, is Step~3(D$_3$) of the main paper.
